\section{Resumen}

Esta tesis estudia el seguimiento de precios en horizonte finito en un creador automático de mercado (AMM) sintético, de producto constante y sin comisiones, frente a un choque en una referencia exógena de un mercado centralizado (CEX). Distingue la mecánica exacta del intercambio de los resultados que dependen de la capacidad de arbitraje, la profundidad, las reglas de los agentes y el orden de eventos. En una piscina aislada, la capacidad $c$ corrige el error logarítmico fuera de banda hasta en $2\log(1+c)$ por tick; para un choque de 0.20 y una banda de 0.003, el plazo ideal es de 20 ticks con $c=0.005$, cuatro con $c=0.03$ y dos con $c=0.10$. El simulador incluye agentes de ruido, tendencia/pánico, arbitraje y provisión de liquidez con retiro por umbral. Se evalúan 13,500 trayectorias principales (300 semillas), 40,500 de sensibilidad (300 semillas nuevas) y una repetición posterior de 9,000 trayectorias con semillas archivadas.

El arbitraje de capacidad alta reduce el área media de exceso de error acumulado de 2.36 sin arbitraje a 0.10, aunque la dirección del efecto se deriva de la regla. Con $c=0.03$, el retorno sostenido de dos ticks comienza antes del plazo de la piscina aislada en el 49.0--84.0\,\% de las trayectorias con profundidad $\kappa=0.5$, el 57.3--83.3\,\% con $\kappa=1$ y el 98.3--100\,\% con $\kappa=2$; todas se recuperan a más tardar en el tick 20 y ninguna orden de las corridas principales alcanza su límite de capital. El umbral de retiro único de E3 se activa en todas las corridas ensayadas, normalmente en el primer tick. Los contrastes suplementarios de respuesta máxima son positivos solo con poca profundidad y órdenes fijas, desaparecen al escalarlas con las reservas y disminuyen al desactivar el pánico. La repetición muestra que las salidas posteriores de E1-alta estático y poca profundidad (90/300 y 53/300 para choques de $-0.05$ y $-0.10$) son episodios únicos asociados al pánico; no hay salidas en las corridas comparables sin pánico. En E3, el 98.2--99.6\,\% de las trayectorias recuperadas vuelve a salir, a menudo repetidamente, con o sin pánico. Fijar la capacidad de arbitraje según las reservas iniciales reduce el contraste del área de exceso posrecuperación de 0.122--0.124 a 0.0036--0.0097. Este control posterior al análisis inicial cubre una profundidad y condiciona el resultado a la recuperación; no demuestra un efecto causal general.

El mapa de dos estados mejora frente a la persistencia en algunas condiciones reservadas, pero no supera al AR(1) escalar, por lo que su segundo estado no demuestra valor predictivo adicional. Un mapa híbrido sensible a la banda reduce cerca del 1.7\,\% el error de reproducción recursiva, pero utiliza la trayectoria futura realizada del CEX y no es un pronóstico en tiempo real. El aporte es una evaluación reproducible y centrada en mecanismos del referente analítico y sus límites. Los resultados son condicionales a un sistema sintético sin comisiones y no tienen validación con datos de mercado; no demuestran efectos de mercado ni requisitos operativos validados. Una reproducción sincronizada de eventos debería registrar ingreso a banda, retorno sostenido, salidas posteriores, órdenes, reservas activas y capacidad de arbitraje disponible. Son observables candidatos, no requisitos validados; su evaluación también requiere comisiones, gas, demoras, costos de operación, límites de capital y competencia entre arbitrajistas realistas.

\textbf{Palabras clave:} creador automático de mercado; arbitraje; seguimiento de precios; retiro de liquidez; simulación basada en agentes; Monte Carlo
